\documentclass[../../main.tex]{subfiles}% 注意这里的文件路径不能用 ./main.tex ,否则用latexmk编译子文件会报错
\graphicspath{{\subfix{./image/}}{\subfix{../../image/}}} % 指定图片引用路径,后续可以直接使用图片文件名
% 如果图片存放在上级目录,则使用 ../../image/ 作为备用路径

% 例如:
% \begin{figure}[H]
% \centering
% \includegraphics[width=0.3\textwidth]{图.png}
% \caption{}
% \label{figure:图}
% \end{figure}
% 注意:上述\label{}一定要放在\caption{}之后,否则引用图片序号会只会显示??.

\begin{document}
% SOURCE p=96; printed=81; object=ch02-notes-title; region=middle
\section{注记与参考文献}\label{section:ch02-notes-references}
% SOURCE p=96; printed=81; object=ch02-notes-lebesgue-history; region=middle
可以说，现代测度与积分理论的历史始于 Lebesgue 于 1902 年发表的博士论文\cite{folland-bib-091}。当然，Lebesgue 的工作建立在其他数学家先前研究的基础上，而且他的一些结果由 Vitali 和 W. H. Young 独立得到。随后十年间，多位数学家广泛发展了 Lebesgue 积分理论，本章大部分结果都是在这一时期首次得到的。特别地，Lebesgue 本人证明了控制收敛定理，并在极限函数 $f$ 可积的情形下由它推出单调收敛定理；当 $\int f=\infty$ 时，后一个定理由 B. Levi 给出。

% SOURCE p=96; printed=81; object=ch02-notes-abstract-history; region=middle-lower
Lebesgue\cite{folland-bib-092}在研究将不定积分概念推广到多变量函数的问题时，考察了 $\mathbb{R}^n$ 上更一般的测度，他称之为“可加函数”。随后，Radon\cite{folland-bib-111}建立了关于现在称为 $\mathbb{R}^n$ 上正则 Borel 测度的积分理论；特别地，$n=1$ 时，这给出了 Lebesgue–Stieltjes 积分。最后，在 1915 年，Fréchet\cite{folland-bib-053}指出，Radon 的许多思想在配备 $\sigma$-代数的集合这一一般背景下同样适用。抽象测度与积分理论由此诞生。它继续发展，直到约 1950 年，才大致形成今天熟悉的样貌。第一部系统的现代专著是 Halmos\cite{folland-bib-061}。

% SOURCE p=96; printed=81; object=ch02-notes-history-sources; region=lower
关于 Lebesgue 积分的前史和早期历史，可参见 Hawkins\cite{folland-bib-070}。关于此后发展的参考文献，可见 Saks\cite{folland-bib-128}。

% SOURCE p=96,97; printed=81,82; object=ch02-notes-daniell; region=bottom/top
本书采取先建立测度、再由测度导出积分的观点。不过，也可以反过来进行，这一方法最初由 Daniell\cite{folland-bib-029}发展。粗略地说，从一个“初等积分”开始：它是定义在适当函数空间上的线性泛函 $I$，满足一些温和的连续性条件，并且在 $f\ge0$ 时有 $I(f)\ge0$，即为正的。例如，$[a,b]$ 上连续函数空间上的 Riemann 积分。Daniell 理论将 $I$ 延拓成定义于更大函数类上的泛函 $\overline I$。在适当假设下，使 $\chi_E$ 属于 $\overline I$ 定义域的集合 $E$ 所组成的族 $\mathscr{M}$ 是一个 $\sigma$-代数，函数 $\mu(E)=\overline I(\chi_E)$ 是 $\mathscr{M}$ 上的测度，而 $\overline I$ 就是关于 $\mu$ 的积分。Daniell 理论的简明介绍见 Royden\cite{folland-bib-121}，全面论述见 Pfeffer\cite{folland-bib-108}；König\cite{folland-bib-086}则给出了稍有不同的方法。

% SOURCE p=97; printed=82; object=ch02-notes-general-integrals-history; region=upper
关于 $\mathbb{R}$ 上的积分，Lebesgue 理论并非终点。部分出于在尽可能一般的条件下建立微积分基本定理的需要（\cref{section:ch03-notes-references}中还会进一步讨论），人们发展了多种积分理论，它们不仅包括 Lebesgue 积分，还包括某些“条件收敛”的积分。也就是说，对于某些满足 $\int f^+=\int f^-=\infty$ 的可测函数 $f:\mathbb{R}\to\mathbb{R}$，通过正值与负值的某种抵消，仍可以合理地定义 $\int f(x)\,\mathrm{d}x$。
\begin{example}\label{example:ch02-conditionally-convergent-integral}
一个标准例子是 $f(x)=x^{-1}\sin x$，见\cref{exercise:ch02-dirichlet-integral}。
\end{example}
最初由 Denjoy 和 Perron 提出的定义此类积分的方法相当复杂。然而，20 世纪 50 年代末，Henstock 和 Kurzweil 独立发现，对经典 Riemann 积分作一种修改就能得到相同结果。

% SOURCE p=97; printed=82; object=ch02-hk-intro; region=middle
有界区间 $[a,b]$ 上的 Henstock–Kurzweil 积分定义如下。
% SOURCE p=97; printed=82; object=ch02-tagged-partition; region=middle
\begin{definition}[带标记的分割]\label{definition:ch02-tagged-partition}
$[a,b]$ 的\textbf{带标记的分割}由有限序列 $\{x_j\}_0^N$ 及另一有限序列 $\{t_j\}_1^N$ 组成，其中 $a=x_0<\cdots<x_N=b$（即\cref{section:ch02-complex-integration}意义下的分割），且 $t_j\in[x_{j-1},x_j]$。
\end{definition}
% SOURCE p=97; printed=82; object=ch02-gauge; region=middle
\begin{definition}[尺度函数]\label{definition:ch02-gauge}
$[a,b]$ 上的\textbf{尺度函数}是一个任意的函数 $\delta:[a,b]\to(0,\infty)$。若 $P$ 为带标记的分割，$\delta$ 为尺度函数，且对所有 $j$，$x_j-x_{j-1}<\delta(t_j)$，则称 $P$ 为 \textbf{$\delta$-细的}。
\end{definition}
% SOURCE p=97; printed=82; object=ch02-fine-partition-existence; region=middle
\begin{proposition}\label{proposition:ch02-fine-partition-existence}
对任意尺度函数 $\delta$，都存在 $[a,b]$ 的 $\delta$-细的带标记分割。
\end{proposition}
\begin{proof}
这容易由 $[a,b]$ 的紧致性推出。
\end{proof}
% SOURCE p=97; printed=82; object=ch02-tagged-riemann-sum; region=middle
\begin{definition}[带标记分割的 Riemann 和]\label{definition:ch02-tagged-riemann-sum}
设 $f$ 为 $[a,b]$ 上的实值函数。若 $P$ 为 $[a,b]$ 的带标记分割，相应的 $f$ 的 \textbf{Riemann 和}为 $\Sigma_Pf=\sum_1^N f(t_j)(x_j-x_{j-1})$。
\end{definition}
% SOURCE p=97; printed=82; object=ch02-henstock-kurzweil-integral; region=middle-lower
\begin{definition}[Henstock–Kurzweil 积分]\label{definition:ch02-henstock-kurzweil-integral}
若存在 $c\in\mathbb{R}$，满足：对每个 $\epsilon>0$，存在尺度函数 $\delta_\epsilon$，使得只要 $P$ 为 $[a,b]$ 的任意 $\delta_\epsilon$-细的带标记分割，就有 $|\Sigma_Pf-c|<\epsilon$，则称函数 $f$ 在 $[a,b]$ 上 \textbf{Henstock–Kurzweil 可积}。此时数 $c$ 唯一，称为 $f$ 的 \textbf{Henstock–Kurzweil 积分}。
\end{definition}
% SOURCE p=97; printed=82; object=ch02-hk-riemann-comparison; region=middle-lower
\begin{remark}
普通 Riemann 积分也可以按完全相同的方式定义，只是仅允许采用常值的尺度函数。
\end{remark}
% SOURCE p=97; printed=82; object=ch02-hk-other-integrals; region=lower
\begin{remark}
Henstock–Kurzweil 积分与 Denjoy 和 Perron 的积分相同。特别地，对非负函数，它与 Lebesgue 积分相同；但其定义域包含许多既取正值又取负值、且不属于 $L^1([a,b])$ 的函数。Henstock–Kurzweil 积分的定义容易推广到无界区间。它还有 $n$ 维版本。
\end{remark}
% SOURCE p=97; printed=82; object=ch02-n-interval-tagged-partition; region=lower
\begin{definition}[$n$ 维区间及其带标记分割]\label{definition:ch02-n-interval-tagged-partition}
将 \textbf{$n$ 维区间}定义为 $n$ 个一维区间的乘积。$n$ 维区间 $I$ 的\textbf{带标记分割}由有限个 $n$ 维区间 $\{I_j\}$ 及每个 $I_j$ 中选取的一个点 $t_j$ 组成，其中 $I_j$ 的内部两两不交，且其并为 $I$。
\end{definition}
积分的定义随后照前面进行。

% SOURCE p=97,98; printed=82,83; object=ch02-hk-evaluation; region=bottom/top
可以提出很有说服力的理由，主张将 Henstock–Kurzweil 积分作为通常向学生讲授的 $\mathbb{R}^n$ 上的积分理论。这不仅因为它更一般，更有力的理由是：它的定义相对简单，入门时不需要测度论。另一方面，它不容易推广到 $\mathbb{R}^n$ 以外的空间；虽然也可以在相当抽象的背景下建立它，但在那里会失去许多原有的简洁性。此外，虽然在某些问题中偶尔会遇到不能由绝对收敛积分通过简单取极限过程得到的条件收敛积分，它们的用途还没有广泛到足以强烈主张让非专门研究者学习的程度。

% SOURCE p=98; printed=83; object=ch02-hk-references; region=upper
无论如何，本书只讨论 Lebesgue 积分及其所属的一般测度与积分理论。希望进一步了解 Henstock–Kurzweil 积分的读者，可以参看 Bartle\cite{folland-bib-013}的简短介绍，以及 McLeod\cite{folland-bib-099}和 Pfeffer\cite{folland-bib-109}的详细论述。此外，Gordon\cite{folland-bib-057}全面介绍了 $[a,b]$ 上的 Denjoy、Perron 和 Henstock–Kurzweil 积分；Henstock\cite{folland-bib-072}则在更抽象的背景下发展了这一理论。

% SOURCE p=98; printed=83; object=ch02-borel-isomorphism; region=middle
\cref{section:ch02-measurable-functions}：
\begin{definition}[Borel 同构]\label{definition:ch02-borel-isomorphism}
两个可测空间 $(X,\mathscr{M})$ 和 $(Y,\mathscr{N})$ 之间的 \textbf{Borel 同构}是一个双射 $f:X\to Y$，使 $f^{-1}$ 是从 $\mathscr{N}$ 到 $\mathscr{M}$ 的双射。
\end{definition}
% SOURCE p=98; printed=83; object=ch02-borel-isomorphism-comment; region=middle
与拓扑空间的同胚概念（见\cref{chapter:ch04-topology}）以及其他各种范畴中的同构概念不同，Borel 同构的用处有限，因为两个空间成为 Borel 同构太容易了。下面这个由 Kuratowski 给出的、该理论中的主要定理清楚地表明了这一点。
% SOURCE p=98; printed=83; object=ch02-kuratowski-borel-isomorphism; region=middle
\begin{theorem}[Kuratowski]\label{theorem:ch02-kuratowski-borel-isomorphism}
设 $(X,\mathscr{M})$ 与完备可分度量空间 $Y$ 的某个 Borel 子集 $E$ Borel 同构，其中 $E$ 配备 $\sigma$-代数 $\{F\in\mathscr{B}_Y:F\subseteq E\}$。则或者 $X$ 可数且 $\mathscr{M}=\mathscr{P}(X)$，或者 $(X,\mathscr{M})$ 与 $(\mathbb{R},\mathscr{B}_{\mathbb{R}})$ Borel 同构。
\end{theorem}
这一结果的证明以及大量关于 Borel 集的其他内容，可见 Srivastava\cite{folland-bib-139}。

% SOURCE p=98; printed=83; object=ch02-baire-hierarchy-intro; region=middle-lower
度量空间上的 Borel 可测函数有一个层级，大致对应于 Borel 集的层级（开集和闭集、$F_\sigma$ 和 $G_\delta$ 等）。
% SOURCE p=98; printed=83; object=ch02-baire-class; region=middle-lower
\begin{definition}[Baire 类]\label{definition:ch02-baire-class}
令 $B_0$ 为所有连续函数组成的空间；对每个可数序数 $\alpha$，递归定义 $B_\alpha$ 如下。若 $\alpha$ 有直接前驱 $\beta$，则 $B_\alpha$ 为 $B_\beta$ 中逐点收敛序列的所有极限组成的集合；否则，$B_\alpha=\bigcup_{\beta<\alpha}B_\beta$。$B_\alpha$ 中的函数称为 \textbf{$\alpha$ 类 Baire 函数}。
\end{definition}
% SOURCE p=98; printed=83; object=ch02-baire-derivative-example; region=lower
\begin{example}\label{example:ch02-baire-derivative}
若 $f$ 在 $\mathbb{R}$ 上处处可微，则 $f'$ 为 $1$ 类 Baire 函数。
\end{example}
% SOURCE p=98; printed=83; object=ch02-notes-exercise-eleven; region=lower
\cref{exercise:ch02-separate-continuity}是 Lebesgue 第一篇已发表论文中的一个结果。相关讨论见 Rudin\cite{folland-bib-123}。

% SOURCE p=98; printed=83; object=ch02-notes-lone-classes; region=lower
\cref{section:ch02-complex-integration}：将单个可测函数与由几乎处处相等关系定义的函数等价类之间的区别模糊处理，通常很方便，而且很少造成严重问题。最常见的需要小心的情形，涉及可测函数与连续函数（例如在 $\mathbb{R}^n$ 上）之间的相互关系，因为一个几乎处处等于连续函数的函数一般并不连续。对此问题的仔细讨论见 Zaanen\cite{folland-bib-165}。

% SOURCE p=98; printed=83; object=ch02-notes-egoroff-lusin; region=bottom
\cref{section:ch02-modes-of-convergence}：Bartle\cite{folland-bib-012}对 Egoroff 定理作了有趣讨论，其中包括几乎一致收敛的一些充要条件。Lusin 定理（\cref{exercise:ch02-lusin}）的一个不依赖 Egoroff 定理的简单证明见 Feldman\cite{folland-bib-043}。我们将在\cref{section:ch07-lp-regular-borel}中证明它的一个更一般形式。

% SOURCE p=99; printed=84; object=ch02-notes-product-history; region=top
\cref{section:ch02-product-measures}：Fubini 和 Tonelli 最初的定理针对平面上的 Lebesgue 测度。抽象乘积测度理论由多人在 20 世纪 30 年代独立发展；这里给出的 $\mu\times\nu$ 构造是 Hahn\cite{folland-bib-060}的构造。还可以在无穷多个测度空间 $\{(X_\alpha,\mathscr{M}_\alpha,\mu_\alpha)\}_{\alpha\in A}$ 的乘积上定义乘积测度，只要除有限多个 $\alpha$ 外均满足 $\mu_\alpha(X_\alpha)=1$；参见 Saeki\cite{folland-bib-127}、Halmos\cite[\S38]{folland-bib-061}或 Hewitt 与 Stromberg\cite[\S22]{folland-bib-076}。我们将在\cref{section:ch07-products-topological-spaces}中给出此结果的一个版本（\cref{theorem:ch07-product-regular-borel}）。

% SOURCE p=99; printed=84; object=ch02-notes-sierpinski; region=upper
Sierpiński\cite{folland-bib-134}利用选择公理，但不用连续统假设，证明了存在 $\mathbb{R}^2$ 的一个非 Lebesgue 可测子集，它与任意直线的交至多包含两个点。这应与\cref{exercise:ch02-nonmeasurable-separately-measurable}比较；该习题也出自 Sierpiński。

% SOURCE p=99; printed=84; object=ch02-generalized-product-measure; region=middle
下面对乘积测度概念的推广在许多情形中很有用。
\Needspace{9\baselineskip}
\begin{proposition}\label{proposition:ch02-generalized-product-measure}
给定可测空间 $(X,\mathscr{M})$、$\sigma$-有限测度空间 $(Y,\mathscr{N},\nu)$，以及 $X$ 上的一族有限测度 $\{\mu_y:y\in Y\}$，满足对每个 $E\in\mathscr{M}$，函数 $y\mapsto\mu_y(E)$ 在 $Y$ 上可测。则可在 $X\times Y$ 上定义一个测度 $\lambda$，使得对 $f\in L^+(X\times Y)$，有
\begin{align*}
\int f\,\mathrm{d}\lambda=\iint f(x,y)\,\mathrm{d}\mu_y(x)\,\mathrm{d}\nu(y).
\end{align*}
\end{proposition}
参见 Johnson\cite{folland-bib-079}。

% SOURCE p=99; printed=84; object=ch02-notes-change-variables; region=middle-lower
\cref{section:ch02-n-dimensional-lebesgue-integral}：\cref{theorem:ch02-differentiable-change}的证明采用 J. Schwartz\cite{folland-bib-131}的方法。对变换 $G$ 作稍弱的假设，也能证明这个定理；参见 Rudin\cite[定理 7.26]{folland-bib-125}。
\end{document}
